% =====================================================================
% Annotation-extraction fixture: ACTUALTEXT
%
%   pdflatex base.tex        (one run is enough; no aux dependencies)
%
% TOPIC: /ActualText -- the document's own declaration of what a run of
% glyphs means, which overrides whatever the font's /ToUnicode says
% about those glyphs.
%
% This is the only place in the corpus where authorial intent is
% RECOVERABLE rather than guessed. Everywhere else an extractor sees a
% glyph and a /ToUnicode entry chosen by the producer's toolchain, and
% cannot tell a typographic ligature from a character the author meant
% to write. /ActualText is the spec's answer to that, and where it is
% present it wins.
%
% Not here: glyph-to-character recovery in the ABSENCE of a declaration
% (tests/text), output sequence (tests/order), page attributes
% (tests/page), the annotation dictionary (tests/inventory,
% tests/comments), dictionaries no producer writes (tests/structure).
%
% HAND-ANNOTATED. /ActualText lives in the content stream, put there by
% the DOCUMENT producer; the annotations over it are ordinary
% highlights. Nothing here needs a generator.
%
% Contains NO annotations. Annotate base.pdf by hand in each program
% following ANNOTATIONS.txt, and save as base-<tool>-<version>.pdf.
%
% Deliberate choices:
%
%   * accsupp for the marked-content wrappers. It is the only widely
%     available way to emit /ActualText from pdflatex. hyperref is still
%     excluded: it would add /Link annotations to /Annots.
%
%   * \pdfglyphtounicode{ff}{FB00} as in tests/text, so that AT1 has a
%     real presentation form to override. Without it AT1's tagged and
%     untagged targets are indistinguishable and the item tests nothing.
%
%   * every \AT call carries the readable string in a comment beside the
%     hex. The hex is UTF-16BE, which is what accsupp's method=hex
%     expects, and is unreadable by eye -- a wrong digit produces a
%     plausible fixture that asserts the wrong thing.
%
%   * uncompressed, so the marked content can be inspected:
%       grep -c ActualText base.pdf          # expect 8
%       grep -cE "/Alt[ (<]" base.pdf        # expect 1
%     A producer that strips marked content on save destroys this whole
%     family silently. Check before trusting a run.
%
%   * every target carries a unique sentinel [[ATn]], at the END of the
%     selection -- EXCEPT AT5A and AT5B, whose selections stop inside a
%     tagged span and so cannot end with one.
% =====================================================================

\documentclass[11pt,a4paper]{article}

\usepackage[T1]{fontenc}
\usepackage[utf8]{inputenc}
\usepackage{lmodern}
\usepackage{accsupp}
\usepackage[margin=3cm]{geometry}

\input glyphtounicode
\pdfgentounicode=1
\pdfglyphtounicode{ff}{FB00}   % see header note

\pdfcompresslevel=0
\pdfobjcompresslevel=0

\pagestyle{plain}
\setlength{\parindent}{0pt}
\setlength{\parskip}{5pt}

% \AT{utf-16be hex}{visible text}  -- declares what the glyphs mean
\newcommand{\AT}[2]{\BeginAccSupp{method=hex,unicode,ActualText=#1}#2\EndAccSupp{}}
% \ATnil{visible text}             -- declares the glyphs mean nothing
\newcommand{\ATnil}[1]{\BeginAccSupp{method=hex,unicode,ActualText=}#1\EndAccSupp{}}
% \ALT{utf-16be hex}{visible text} -- an ALTERNATE DESCRIPTION, not a
%                                     replacement: must never be substituted
\newcommand{\ALT}[2]{\BeginAccSupp{method=hex,unicode,Alt=#1}#2\EndAccSupp{}}

\begin{document}

\section*{AT1 --- a declaration that decomposes a ligature}

The tagged target and its untagged control render identically. Only the
declaration tells them apart, and only when ligatures are kept.

the tagged \AT{0074006100720069006600660073}{tariffs} of the schedule [[AT1A]]
% ActualText = tariffs

the untagged tariffs of the schedule [[AT1B]]

\section*{AT2 --- a declaration that asserts a presentation form}

Here the glyphs decompose on their own and the document declares the
presentation form anyway. This is the case a fold cannot be allowed to
undo: the author asked for U+FB01.

a declared \AT{FB01007800650064}{fixed} point in the mapping [[AT2A]]
% ActualText = U+FB01 xed

\section*{AT3 --- declared text of a different length}

One glyph standing for many characters, and many glyphs standing for
one. Both directions, because a length assumption breaks only one way.

the \AT{00530065006300740069006F006E}{\S} on rounding conventions [[AT3A]]
% ActualText = Section

the \AT{2122}{(TM)} mark on the product name [[AT3B]]
% ActualText = U+2122

\section*{AT4 --- a hyphen declared to mean nothing}

The break below is inside one word and its hyphen carries an EMPTY
declaration, so the join is not a guess. Compare tests/text TX4, where
nothing in the PDF distinguishes this from a compound hyphen.

\parbox{0.9\textwidth}{\raggedright
the surrounding infra\ATnil{-}\linebreak structure of the market [[AT4A]]%
}

\section*{AT5 --- a selection that covers part of a tagged span}

A declaration applies to its whole span. A quad that covers only part
of one cannot be told which part of the declared text corresponds.

[[AT5A]] see \AT{00460069006700750072006500200033}{the third figure} here
% ActualText = Figure 3

[[AT5B]] consult \AT{00460069006700750072006500200033}{the third figure} here
% ActualText = Figure 3

\section*{AT6 --- an alternate description, not a replacement}

/Alt describes; /ActualText replaces. Substituting /Alt into the
extracted text is a defect, not a nicety.

the \ALT{0041004C00540020004D0055005300540020004E004F00540020004100500050004500410052}{described}
% Alt = ALT MUST NOT APPEAR
phrase in this sentence [[AT6A]]

\section*{AT7 --- a tagged span across a line break}

One declaration, two quads. The declared text must appear once.

\parbox{0.42\textwidth}{\raggedright
the \AT{00460069006700750072006500200033}{third figure\linebreak of the set}
% ActualText = Figure 3, and the span itself is broken across two lines
is cited here [[AT7A]]%
}

\end{document}
